\section{Superinteligencia: un agent de aprendizaje continuo}

\subsection{La redefinición del AGI}

\begin{importantbox}{De «lo sabe todo» a «puede aprenderlo todo»}
Ilya propone una redefinición de la superinteligencia: no una «finished mind which knows how to do every single job», sino una \textbf{mind capaz de aprender rápidamente cualquier trabajo}. Como un «adolescente de 15 años superinteligente»: no sabe nada, pero aprende cualquier cosa con extrema rapidez.
\end{importantbox}

El proceso de despliegue incluirá etapas de aprendizaje y prueba y error: la IA se despliega en distintos puestos de la economía y aprende de manera continua en el trabajo (on-the-job continual learning), igual que un empleado humano recién contratado.

\subsection{Dos caminos hacia la superinteligencia}

\begin{enumerate}
    \item \textbf{Automejora recursiva}: si este algoritmo de aprendizaje eficiente también se vuelve superhumano en tareas de investigación de aprendizaje automático, el algoritmo mismo se mejora continuamente, formando un ciclo de aceleración
    \item \textbf{Fusión de conocimiento}: incluso sin automejora recursiva, distintas instancias de un mismo modelo aprenden en distintos puestos de la economía y después \textbf{fusionan lo aprendido} (amalgamating the learnings); los humanos no pueden fusionar sus mentes, pero la IA sí puede: esto ya constituye por sí mismo una superinteligencia funcional
\end{enumerate}

\subsection{Crecimiento económico acelerado}

Ilya considera que un despliegue amplio podría traer un \textbf{crecimiento económico acelerado}. Las distintas reglas de cada país llevarán a distintas velocidades de crecimiento: los países con reglas más favorables crecerán más rápido. Pero la velocidad concreta es difícil de predecir, porque cambiar el mundo físico sigue tomando tiempo.

\subsection{Resumen de la sección}

La superinteligencia no es un oracle omnisciente, sino un agent que aprende con extrema rapidez. Mediante el aprendizaje continuo y la fusión de conocimiento, incluso sin automejora recursiva, se puede alcanzar una superinteligencia funcional.


